% =========================================================================
\section{Circuito 3: amplificador inversor generalizado y filtro paso bajo}\label{sec:pl-filtro}
% =========================================================================
\begin{prelab}[Inversor generalizado y filtro]\label{pl:filtro}
Sea el circuito de la \figpanel{fig:generalizado}{a}, en el que las impedancias $Z_1$ y $Z_2$ pueden
ser no resistivas y, por lo tanto, tener dependencia en frecuencia: obténgase la ganancia en voltaje
del circuito. Esta ganancia será función de la frecuencia salvo que ambas impedancias sean puramente
resistivas. Considérese el caso de que $Z_1=R$ sea real y $Z_2$ una composición en paralelo de una
resistencia $R_F$ y un condensador $C$ (\figpanel{fig:generalizado}{b}). Demuéstrese que la dependencia
de la función de transferencia con la frecuencia $f$ es la de un filtro paso bajo: para frecuencias
suficientemente bajas la ganancia en tensión es constante, mientras que para frecuencias
suficientemente altas disminuye.
\end{prelab}

% Figura: (a) amplificador inversor generalizado; (b) caso Z1=R, Z2=RF||C (filtro paso bajo).
\begin{figure}[ht]
\centering
\begin{minipage}[b]{0.46\textwidth}\centering
\begin{circuitikz}[circuito, scale=0.88, transform shape]
  \draw (0,0) node[op amp, ao] (A) {};
  \draw (A.-) -- ++(-0.5,0) coordinate (N);
  \draw (N) to[generic, l_=$Z_1$, i<^=$I_1$] ++(-2.2,0) to[short, -o] ++(-0.3,0) node[left, tens] {$\Vin$};
  \draw (A.out) -- ++(0.6,0) coordinate (O);
  \draw (N) -- ++(0,1.4) coordinate (T) to[generic, l=$Z_2$, i>^=$I_2$] (T -| O) -- (O);
  \draw (O) to[short, *-o] ++(0.6,0) node[right, tens] {$\Vout$};
  \draw (A.+) -- ++(-0.5,0) -- ++(0,-0.4) node[ground] {};
  \fill (N) circle (1.6pt);
\end{circuitikz}\\[2pt]
{\small (a)}
\end{minipage}\hfill
\begin{minipage}[b]{0.50\textwidth}\centering
\begin{circuitikz}[circuito, scale=0.88, transform shape]
  \draw (0,0) node[op amp, ao] (A) {};
  \draw (A.-) -- ++(-0.5,0) coordinate (N);
  \draw (N) to[R, l_=$R$] ++(-2.2,0) to[short, -o] ++(-0.3,0) node[left, tens] {$\Vin$};
  \draw (A.out) -- ++(0.6,0) coordinate (O);
  \draw (N) -- ++(0,1.2) coordinate (T1) to[C, l_=$C$] (T1 -| O) -- (O);
  \draw (T1) -- ++(0,1.0) coordinate (T2) to[R, l=$R_F$] (T2 -| O) -- (T1 -| O);
  \fill (T1) circle (1.6pt); \fill (T1 -| O) circle (1.6pt);
  \draw (O) to[short, *-o] ++(0.6,0) node[right, tens] {$\Vout$};
  \draw (A.+) -- ++(-0.5,0) -- ++(0,-0.4) node[ground] {};
  \fill (N) circle (1.6pt);
\end{circuitikz}\\[2pt]
{\small (b)}
\end{minipage}
\caption{(a) Amplificador inversor generalizado (figura~4a del guion~\cite{GUION-S1}): $Z_1(s)$ y $Z_2(s)$
son impedancias arbitrarias. (b) Caso $Z_1=R$, $Z_2=R_F\parallel C$ (figura~5a del guion), que es un
filtro paso bajo activo y, para $R_F\to\infty$, un integrador.}
\label{fig:generalizado}
\end{figure}


% -------------------------------------------------------------------------
\subsection{Ganancia del inversor generalizado}\label{sec:pl-generalizado}
% -------------------------------------------------------------------------
\begin{resolucion}
Trabajamos en el dominio de Laplace con condiciones iniciales nulas: cada elemento cumple
$V(s)=Z(s)\,I(s)$, \cref{eq:Zs}, y Kirchhoff, (R1) y (R2) se mantienen por linealidad. El cálculo es
el del inversor con $R_k\to Z_k(s)$.

\textbf{1. Tierra virtual} (R2):
\begin{equation}\label{eq:gen-vm}
  V_-(s)=V_+(s)=0 .
\end{equation}

\textbf{2. Ohm en $Z_1$ y Kirchhoff en el nudo inversor} ($I_-=0$, R1):
\begin{equation}\label{eq:gen-I}
  I_2(s)=I_1(s)=\frac{\Vin(s)}{Z_1(s)} .
\end{equation}

\textbf{3. Ohm en $Z_2$:}
\begin{equation}\label{eq:gen-vout}
  \Vout(s)=V_-(s)-Z_2(s)\,I_2(s)=-\frac{Z_2(s)}{Z_1(s)}\,\Vin(s) .
\end{equation}

\textbf{Resultado.} Con $s=\jj\omega$ se obtiene la ganancia en régimen sinusoidal (\cref{prop:sinusoide}):
\begin{equation}\label{eq:gen-H}
  \resultado{H(s)=\frac{\Vout(s)}{\Vin(s)}=-\frac{Z_2(s)}{Z_1(s)}},\qquad
  A_V(\jj\omega)=-\frac{Z_2(\jj\omega)}{Z_1(\jj\omega)} .
\end{equation}
Si $Z_1=R_1$ y $Z_2=R_2$ se recupera \cref{eq:inv-Av}, independiente de $\omega$; si alguna impedancia
contiene un condensador, $Z_C=1/(\jj\omega C)$ introduce la dependencia en frecuencia.
\end{resolucion}

% -------------------------------------------------------------------------
\subsection{Caso $Z_1=R$, $Z_2=R_F\parallel C$: filtro paso bajo}\label{sec:pl-pasobajo}
% -------------------------------------------------------------------------
\begin{resolucion}
\textbf{1. Impedancias} (paralelo, \cref{eq:RC-paralelo}):
\begin{equation}\label{eq:pb-Z}
  Z_1=R,\qquad
  Z_2(s)=\frac{R_F\cdot\frac{1}{sC}}{R_F+\frac{1}{sC}}=\frac{R_F}{1+sR_FC} .
\end{equation}

\textbf{2. Función de transferencia}, sustituyendo \cref{eq:pb-Z} en \cref{eq:gen-H}:
\begin{equation}\label{eq:pb-H}
  \resultado{H(s)=-\frac{R_F}{R}\,\frac{1}{1+sR_FC}}
\end{equation}
Es una red de constante temporal única paso bajo, \cref{eq:STC}, con
\begin{equation}\label{eq:pb-K}
  K=-\frac{R_F}{R},\qquad \tau=R_FC,\qquad f_c=\frac{1}{2\pi R_FC} .
\end{equation}

\textbf{3. Respuesta en frecuencia} ($s=\jj2\pi f$, \cref{eq:STC-modulo}):
\begin{equation}\label{eq:pb-modulo}
  \abs{A_V(f)}=\frac{R_F/R}{\sqrt{1+(f/f_c)^{2}}},\qquad
  \arg A_V(f)=180^{\circ}-\arctg\frac{f}{f_c} .
\end{equation}

\textbf{4. Límites.}
\begin{align}
  f\ll f_c:&\quad \abs{A_V}\simeq\frac{R_F}{R}
     &&\text{($C$ abierto: constante)},\label{eq:pb-baja}\\
  f\gg f_c:&\quad \abs{A_V}\simeq\frac{R_F}{R}\,\frac{f_c}{f}=\frac{1}{2\pi fRC}
     &&\text{($C$ cortocircuita a $R_F$: decrece, $-20\dB$/década)},\label{eq:pb-alta}
\end{align}
y en $f=f_c$, $\abs{A_V}=(R_F/R)/\sqrt2$ ($-3\dB$). Es la dependencia de la figura~4b del guion
(\cref{fig:bode-filtro}): el circuito es un \emph{filtro paso bajo de primer orden}.
\end{resolucion}

% Figura: Bode del filtro paso bajo activo con R=1k, RF=100k, C=1,59 nF (fc=1 kHz).
\begin{figure}[ht]
\centering
\begin{tikzpicture}
  \begin{semilogxaxis}[tema, name=mod, width=0.86\textwidth, height=5.6cm,
      xmin=10, xmax=1e6, ymin=-22, ymax=48, ytick={-20,0,20,37,40},
      xlabel={$f$ (Hz)}, ylabel={$\abs{A_V}$ (dB)}, log basis x=10,
      title style={font=\small}, title={(a) Módulo}, legend pos=south west]
    \addplot[azulacero, line width=1.1pt, domain=10:1e6, samples=300]
      {40-10*log10(1+(x/1001)^2)};
    \addlegendentry{{set 1 ($R_F=100\,\mathrm{k}\Omega$)}}
    \addplot[rojoteo, line width=0.9pt, densely dashed, domain=10:1e6, samples=2]
      {20*log10(100097/x)};
    \addlegendentry{{set 2 ($R_F\to\infty$, integrador)}}
    \addplot[asintota, domain=10:1001, samples=2] {40};
    \addplot[punteado] coordinates {(1001,-22) (1001,37)};
    \addplot[punteado] coordinates {(100097,-22) (100097,0)};
    \addplot[only marks, mark=*, mark size=1.6pt, azulacero] coordinates {(1001,37)};
    \addplot[only marks, mark=*, mark size=1.6pt, rojoteo] coordinates {(100097,0)};
    \node[font=\scriptsize, anchor=south west] at (axis cs:1150,37.5) {$-3\dB$ en $f_c\approx1\,\mathrm{kHz}$};
    \node[font=\scriptsize, anchor=south west] at (axis cs:1.1e5,0.5) {$f_1\approx100\,\mathrm{kHz}$};
    \node[font=\scriptsize, anchor=south west, text=grispizarra, rotate=-24] at (axis cs:6000,22) {$-20\dB$/década};
  \end{semilogxaxis}
  \begin{semilogxaxis}[tema, at={(mod.south)}, anchor=north, yshift=-1.4cm,
      width=0.86\textwidth, height=4.2cm,
      xmin=10, xmax=1e6, ymin=80, ymax=190, ytick={90,135,180},
      xlabel={$f$ (Hz)}, ylabel={$\arg A_V$ (grados)}, log basis x=10,
      title style={font=\small}, title={(b) Fase}]
    \addplot[azulacero, line width=1.1pt, domain=10:1e6, samples=300] {180-atan(x/1001)};
    \addplot[rojoteo, line width=0.9pt, densely dashed, domain=10:1e6, samples=2] {90};
    \addplot[only marks, mark=*, mark size=1.6pt, azulacero] coordinates {(1001,135)};
    \node[font=\scriptsize, anchor=west] at (axis cs:1200,137) {$135^{\circ}$};
  \end{semilogxaxis}
\end{tikzpicture}
\caption{Diagrama de Bode del filtro de la \figpanel{fig:generalizado}{b} con los valores del set~1,
según \cref{eq:pb-modulo}: meseta de $40\dB$ (ganancia $100$), corte a $-3\dB$ en $f_c\approx1\,\mathrm{kHz}$ y
caída de $-20\dB$/década. En trazo discontinuo, el integrador ideal del set~2, que coincide con la
asíntota de alta frecuencia y corta $0\dB$ en $f_1=1/(2\pi RC)\approx100\,\mathrm{kHz}$. La fase pasa de
$180^{\circ}$ (inversión) a $90^{\circ}$ (integración), con $135^{\circ}$ en $f_c$.}
\label{fig:bode-filtro}
\end{figure}

